\chapter{\rm\bfseries Experimental Design and Procedure}
\label{ch:experiments}

The empirical campaign was organised into two sequential phases on a
single experimental apparatus. Phase~1 established the contact-force
versus press-depth relationship for each gripper variant at nominal
zero-tilt orientation. Phase~2 then perturbed the same measurement with
$\pm 2.5^{\circ}$ tilts about the two orthogonal Cartesian axes.
Together the two phases define a three-dimensional
(depth, tilt-axis, tilt-sign) cell grid. The per-cell measurements of
that grid constitute the data corpus reported in
Chapter~\ref{ch:results}.

\section{Study Design}
\label{sec:study-design}

The gripper roster comprised six variants, organised into two families.
\begin{itemize}
  \item \textbf{Finray family}, four variants. The family comprises
    \texttt{finray\_18}; \texttt{finray\_18\_model2}, a second print of
    the same CAD; \texttt{finray\_26}, with a larger blade; and
    \texttt{finray\_26\_5deg}, which is \texttt{finray\_26} on a
    $5^{\circ}$-inclined seat.
  \item \textbf{Parallel-jaw family}, two variants. These are
    \texttt{parallel\_jaw\_stock}, the bare rigid jaws, and
    \texttt{parallel\_jaw\_TPU}, the same jaws with a $2$\,mm TPU $90$A
    compliant pad.
\end{itemize}

The same fabric specimen (single-layer cotton) rested on the same
acrylic table surface throughout the campaign. The Franka FR3
arm held the gripper in a fixed base orientation $R_0 = R_x(180^{\circ})$
for Phase~1. For Phase~2, signed rotations
$R_{\theta} = R_x(180^{\circ})\, R_y(\theta_y)\, R_x(\theta_x)$
were applied, with $(\theta_x, \theta_y) \in
\{(\pm 2.5^{\circ}, 0), (0, \pm 2.5^{\circ})\}$. These are the runtime
arm tilts of Chapter~\ref{ch:methods}, illustrated in
Figure~\ref{fig:runtime-tilt}. They are distinct from the seat angle
$\alpha$ of Figure~\ref{fig:seat-mount-params}.

The perturbation magnitudes were chosen to represent the placement
uncertainty of a deployment that cannot afford precise calibration. A
low-accuracy or coarsely mounted robot arm may plausibly carry an
orientation error of a few degrees and a surface-height error of
several millimetres. The $\pm 2.5^{\circ}$ runtime tilt stands in for
that orientation error. The depth sweep spans the corresponding
height-error budget, and its $10$\,mm calibrated analysis window is
defined in Chapter~\ref{ch:methods}. Therefore, a gripper that
tolerates the full tested range also tolerates a deployment in which
neither the mounting angle nor the surface height is precisely known.

Table~\ref{tab:test-conditions} lists the test conditions common
to both phases.

\begin{table}[ht]
\centering
\caption[Test conditions]{Test conditions of the experimental campaign.
Symbols are defined in the Notation section; $\theta_x$, $\theta_y$ are
illustrated in Figure~\ref{fig:runtime-tilt}.}
\label{tab:test-conditions}
\small
\begin{tabular}{lll}
\toprule
Variable & Definition & Values / precision \\
\midrule
Press depth $d$ & Fingertip driven below the & rigid: $-1$ to $+4$\,mm at \\
                & detected surface           & \quad $0.25$--$0.5$\,mm steps \\
                &                            & finray: $-2$ to $+16$\,mm at \\
                &                            & \quad $0.5$--$1.0$\,mm steps \\
Tilt $\theta_x$ & Arm leaned about the symmetric axis & $-2.5^{\circ}$, $0^{\circ}$, $+2.5^{\circ}$ \\
Tilt $\theta_y$ & Arm leaned along the gripping direction & $-2.5^{\circ}$, $0^{\circ}$, $+2.5^{\circ}$ \\
Surface reference & Force touch-off at 3\,N threshold & repeatable to $\pm 0.1$\,mm \\
Repeats & Trials per (variant, depth, tilt) cell & 3 \\
\bottomrule
\end{tabular}
\end{table}

\section{Depth Grid}
\label{sec:depth-grid}

The depth axis is defined relative to the per-session detected cloth
surface via the \texttt{ground\_truth\_finder} touch-off procedure
described in Chapter~\ref{ch:methods}. The depth grid was chosen
adaptively for each gripper family.
\begin{itemize}
  \item For the rigid parallel jaws, depths were spaced at $0.25$--$0.5$\,mm
        intervals over $-1$ to $+4$\,mm, close to the narrow
        force-limited working range.
  \item For the finrays, depths were spaced at $0.5$--$1.0$\,mm intervals
        over $-2$ to $+16$\,mm, covering the four regimes
        (boundary, knee, valley, recovery) identified in
        Chapter~\ref{ch:results}.
\end{itemize}

Three trials were collected per (variant, depth, tilt) cell. Boundary
depths were routinely retested when the within-cell standard deviation
exceeded $1$\,N, up to a maximum of two additional trials per
(depth, tilt) combination.

\section{Tilt-Aware Pose Geometry: Verification}
\label{sec:tilt-verify}

Before reporting force results, the implementation of the tilt-aware
tool-centre-point (TCP) tracking equation was verified empirically.
Table~\ref{tab:tilt-geometry} shows the measured flange XY shift during
the stationary grasp sample versus the geometric prediction for
\texttt{finray\_26} at four representative (depth, tilt) cells.

\begin{table}[ht]
\centering
\caption[Tilt geometry verification]{Measured versus predicted flange-XY
offset during the stationary grasp sample at four cells. Predicted
offset uses $z_{\text{TCP}} = 94.911$\,mm.}
\label{tab:tilt-geometry}
\footnotesize
\begin{tabular}{lccccc}
\toprule
Cell & Tilt cmd & Predicted $\Delta \mathbf{p}_\mathcal{F}$ (mm) &
       Measured $\Delta \mathbf{p}_\mathcal{F}$ (mm) & Residual & Mean $F_z$ \\
\midrule
$-1.0$\,mm   & $+2.5^{\circ}$ y & $(+4.14, +0.00)$ & $(+4.35, -0.03)$ & $0.21$\,mm & $-8.00$\,N \\
$-1.0$\,mm   & $-2.5^{\circ}$ y & $(-4.14, +0.00)$ & $(-4.02, -0.04)$ & $0.13$\,mm & $-3.39$\,N \\
$+10.0$\,mm  & $+2.5^{\circ}$ y & $(+4.14, +0.00)$ & $(+5.35, -0.28)$ & $1.24$\,mm & $-36.26$\,N \\
$+10.0$\,mm  & $-2.5^{\circ}$ y & $(-4.14, +0.00)$ & $(-2.93, -0.29)$ & $1.24$\,mm & $-35.95$\,N \\
\bottomrule
\end{tabular}
\end{table}

At boundary depth, the residual was $\leq 0.21$\,mm, which is
approximately the lateral positioning noise of the arm under light
load. At safe depth, both positive and negative tilts exhibited a
$+1.2$\,mm common-mode offset in the same $+x$ direction, regardless of
tilt sign. This common-mode offset is interpreted as evidence of finger
compliance under load. The $\sim 36$\,N axial press deflects the soft
finray blade laterally by approximately $1.2$\,mm in a structurally
preferred direction. The differential offset (between $+y$ and $-y$
tilts) remains correctly predicted by the geometric model in both depth
conditions.

This verification confirms that the control-frame implementation of the
tilt-aware TCP tracker is correct. Any asymmetries observed in the
contact force therefore reflect physical gripper behaviour instead of
control-system errors.

\section{Rotation Tracking Verification}
\label{sec:rotation-verify}

The tilt conditions assume that the Franka arm realises the commanded
orientation faithfully. The 30\,Hz session log records actual and
commanded RPY at every sample, giving a direct measurement of the
orientation tracking error
\begin{equation}
  \delta(t) = \| \mathrm{RPY}_{\text{actual}}(t) -
                 \mathrm{RPY}_{\text{cmd}}(t) \|_2.
\end{equation}
Across five representative sessions, the mean steady-state $\delta$
was $0.05^{\circ}$ with a worst-case $0.13^{\circ}$. Transient
excursions up to $\sim 5^{\circ}$ occur exclusively within the
\path{orientation_lock} stage at session start, before any descent
or grasp. Such excursions are filtered out of trial-level analyses.

A $2.5^{\circ}$ commanded tilt is therefore $50\times$ the noise floor
of the arm's rotation tracking. The observed force differences cannot
be attributed to arm rotation error.

\section{Trial Protocol and Outcome Classification}
\label{sec:trial-protocol}

Each trial followed the fixed stage sequence detailed in
Section~\ref{sec:trial-procedure}. The commanded pose was set above the
target, and the arm descended to the pre-contact height. The gripper
then contacted the cloth and pressed to the target depth, closed with
its default close command, and lifted by a fixed clearance. The
operator's held/not-held verdict was recorded at the end of the trial.
A stationary-sample window of $1$\,s (30 samples) during the
grasped-and-pressed stage provided the force measurement.

Trial outcomes were classified into three categories.
\begin{description}
  \item[\textbf{Success}:] the trial executed to completion, contact
        force was sampled cleanly during the stationary window, and the
        cloth remained gripped after the lift stage.
  \item[\textbf{No-grasp}:] the trial executed to completion and
        contact force was sampled cleanly during the stationary window,
        but the cloth slipped from the fingers during or after the lift
        stage. The force measurement remains valid as an
        \emph{engagement-force} characterisation.
  \item[\textbf{Failed trial}:] the trial did not execute to completion.
        The controller either aborted with an excess-force safety event
        or the operator interrupted before the stationary window
        completed. The measurement is invalid and the trial is excluded
        from analyses.
\end{description}

Section~\ref{sec:failure-mode-discussion} discusses why this three-way
taxonomy was adopted in place of a binary success/failure
classification.

\section{Experiment 3: Inclined-Surface Demonstration (Pilot)}
\label{sec:experiment3}

Experiments 1 and 2 characterise the contact phase on a flat, levelled
surface, and each variant is calibrated against that surface.
Experiment 3 tests the deployment consequence that those measurements
predict. In this experiment the surface geometry departs from what the
calibration assumed, while the calibration, the control policy and the
sensing are left unchanged. The question is whether the compliant
finger keeps grasping under that departure where the rigid jaw fails.

The planned protocol mirrors Experiments 1 and 2 in every respect except
the surface and the outcome metric. The acrylic test surface is mounted
on a wedge, \texttt{SURFACE\_RISER}, whose CAD model specifies a
top-surface incline of $\beta = 4.15^{\circ}$. Measured from the model
geometry, the top face normal is $4.1467^{\circ}$ from vertical. The
slope is aligned with the gripping direction, and the as-built fixture
will be verified against the nominal value before the session. Each
variant is then calibrated \emph{blind} by the standard force touch-off
of Section~\ref{sec:tcp-calib}, exactly as if the surface were flat and
level. No tilt compensation is applied, and the incline is treated as
unknown to the calibration. Each trial presses to the calibrated
touch-off depth plus $0.5$\,mm, which is the shallowest depth at which
every variant grasped successfully on the flat surface. The outcome is
recorded as a binary grasp verdict under the outcome taxonomy of
Section~\ref{sec:trial-protocol}. Contact force is logged, although it
is not a reported metric for this experiment.

The pilot covers all six gripper variants of
Table~\ref{tab:variants}, at three attempts each, giving $18$ trials
in total.
From the flat-surface working ranges of Chapter~\ref{ch:results}, the
registered prediction is that all four finray variants grasp
successfully at every attempt, while both rigid variants fail or are
marginal. A $4.15^{\circ}$ surface-relative misalignment exceeds the
largest runtime tilt at which the rigid jaws retained a usable window
($2.5^{\circ}$). Under that tilt the window had already narrowed to
between $0.5$ and $2.5$\,mm. In addition, one
\texttt{parallel\_jaw\_stock} tilt condition was reduced to a single
successful depth, and two produced no success at all
(Section~\ref{sec:pj-vs-finray-results}).

\begin{figure}[ht]
\centering
\includegraphics[width=0.95\linewidth]{exp3_wedge.jpg}
\caption[Experiment 3 wedge fixture]{The \texttt{SURFACE\_RISER} wedge
(3D-printed, white) mounted under the acrylic test surface, on the
optical breadboard used for Experiments 1 and 2. The arrow marks the
gripper's local $y$ axis (the gripping direction), along which the
wedge slope is aligned.}
\label{fig:exp3-wedge}
\end{figure}

\TODO{Experiment 3 is a placeholder: the protocol above is specified
(Bill 0009, pending enactment) and the wedge fixture is built (Figure~\ref{fig:exp3-wedge}), but the pilot
has not yet been run. Insert the outcome grid (variant $\times$ attempt,
success/no-grasp/abort), the as-built incline angle if it differs from
the CAD nominal $4.15^{\circ}$, and a short results paragraph once the
session is complete. If the pilot is not run before submission, delete
this section and retain the inclined-surface demonstration as future work
in Chapter~\ref{ch:conclusion}.}
